\section{A second one-pole extension over prime fields}
\label{app:six-word-one-pole}

The following is a fixed-size positive construction.  It does not assert an
unbounded list, a superlinear line-exception count, or a benchmark improvement.

\begin{theorem}\label{thm:six-word-one-pole}
There are a number field $K$, a set $\Omega\subset K$ of sixty distinct
points, a word $w\colon\Omega\to K$, and six distinct polynomials of degree
at most fourteen, each agreeing with $w$ on at least thirty points.
The same statement holds over $\mathbb F_p$ for arbitrarily large primes $p$.
Thus the code dimension is fifteen, exactly one quarter of the length.
\end{theorem}

\begin{proof}
We first construct a five-polynomial word of length twenty-eight together
with a proper one-pole rational function having fourteen agreements.

\paragraph{Four quadratics and the first rational witness.}
Choose nonzero parameters $a_1,\ldots,a_4$ with distinct squares and with
all twelve signed pair products distinct.  Put
\[
 H_i(X)=X^2/a_i^2+a_i^2,\qquad
 \Omega_0=\{\pm a_i a_j:i<j\}.
\]
At either signed node for the pair $ij$, give the received word the value
$a_i^2+a_j^2$.  Exactly $H_i,H_j$ agree there, and every $H_i$ has six
agreements.  Write $e_j$ for the elementary symmetric functions of the
four parameters and set
\begin{align*}
 D_1(X)&=e_1e_4-e_3X,\\
 N_1(X)&=e_1X^3+(e_3-e_1e_2)X^2
                 +(e_2e_3-e_1e_4)X-e_3e_4.
\end{align*}
The coefficient identity
\[
 N_1-D_1H_i
 =\frac{\prod_{j\ne i}(a_i+a_j)}{a_i^2}
       \prod_{j\ne i}(X-a_i a_j)
\]
follows by using
$a_i^4-e_1a_i^3+e_2a_i^2-e_3a_i+e_4=0$.
Assume $e_3\ne0$ and that
$\alpha=e_1e_4/e_3$ lies outside $\Omega_0$.
Then $N_1/D_1$ is proper and agrees at the six positive pair nodes.
Indeed, a common zero of $N_1,D_1$ would, by the identity, lie among the
positive incident edges for every vertex, which is impossible.

Choose $c$ outside $\Omega_0\cup\{\alpha\}$ and put $y=T^2+c$.
On the twenty-six-point core
\[
 \mathcal C=\{T:T^2+c\in\Omega_0\cup\{\alpha\}\}
\]
use the word
\[
 W(T)=(y-\alpha)(a_i^2+a_j^2)\quad(y=\pm a_i a_j),
 \qquad W(T)=0\quad(y=\alpha).
\]
The five polynomials
\[
 G_i(T)=(y-\alpha)H_i(y)\quad(1\le i\le4),
 \qquad G_5(T)=-N_1(y)/e_3
\]
have degree at most six.  The first four have fourteen core agreements;
$G_5$ has twelve, over the six positive pair nodes.  They are distinct,
since equality with $G_5$ would make the first rational witness polynomial.
All fibers here have two distinct points.

\paragraph{An exact finite-field certificate for a second rational witness.}
The following data are in $\mathbb F_{97}$:
\[
 (a_1,a_2,a_3,a_4)=(48,1,31,32),\quad
 \alpha=53,\quad c=0,\quad b=81,
\]
\begin{align*}
 N(T)&=34+54T+93T^2+84T^4+3T^5+56T^6+70T^7,\\
 G_5(T)&=86+17T^4+93T^6.
\end{align*}
The proper function $N(T)/(T-b)$ agrees at the following twelve core
points.  The first row is $T$, the second is $y=T^2+c$, and the third is
$W(T)$:
\[
\begin{array}{c|rrrrrrrrrrrr}
 T&76&57&90&79&8&88&93&39&56&68&33&47\\
 y&53&48&49&33&64&81&16&66&32&65&22&75\\
 W&0&18&92&78&25&64&54&90&9&78&60&20
\end{array}
\]
These are one pole-fiber point and one point above every old base node
except the positive edge $a_2a_3=31$.  Direct substitution verifies the
agreements.  The pole is outside the entire core and $N(81)=48\ne0$.
Moreover, exactly five displayed points are agreements of $G_5$, and
\[
 N(T)-(T-81)G_5(T)
 =\prod_{t\in\{57,79,88,56,33\}}(T-t)\,
       (74T^2+2T+8).
\]
The remaining quadratic has discriminant $61\ne0$ and roots $24,52$,
both outside the core and the pole.  Giving these two fresh coordinates
the values of $G_5$ therefore produces the desired length-twenty-eight
word over $\mathbb F_{97}$.  To transfer this construction to
characteristic zero, it remains essential to verify smooth lifting;
the finite-field example alone would not suffice.

\paragraph{Polynomial equations and a nonsingular lifting certificate.}
Write
\[
 N(T)=E(y)+T O(y),\qquad
 E(y)=\sum_{j=0}^3 u_jy^j,\quad O(y)=\sum_{j=0}^3 v_jy^j.
\]
Let $\mathcal Y$ consist of $\alpha$ and all signed pair products except
$+a_2a_3$.  For $y=\alpha$ put $w_y=0$; for
$y=\pm a_i a_j$ put $w_y=(y-\alpha)(a_i^2+a_j^2)$.
Consider the following thirteen polynomial equations in fifteen variables,
ordered as $(a_1,\ldots,a_4,\alpha,c,b,u_0,\ldots,u_3,v_0,\ldots,v_3)$:
\begin{equation}\label{eq:six-word-norm-system}
 e_3\alpha-e_1e_4=0,\qquad
 (E(y)+b w_y)^2-(y-c)(O(y)-w_y)^2=0
       \quad(y\in\mathcal Y).
\end{equation}
Their solution at the displayed point is
\[
 (48,1,31,32,53,0,81,34,93,84,56,54,0,3,70).
\]
The $13\times13$ Jacobian minor on the first thirteen columns is nonzero
modulo $97$.  Exact modular differentiation and elimination certify rank
thirteen; the reproducible certificate is
\texttt{transversal\_search/independent\_lift.py} with its accompanying JSON.
At all twelve nodes $O(y)-w_y\ne0$, so the selected root is recovered by
\[
 T=-\frac{E(y)+b w_y}{O(y)-w_y}.
\]
Consequently these norm equations do not introduce a spurious square-root
choice near the certified point.

There is also an independently checked pointed certificate: use variables
\[
 (a_1,\ldots,a_4,\alpha,c,b,N_0,\ldots,N_7,T_0,\ldots,T_{11})
\]
and equations $e_3\alpha-e_1e_4=0$, followed, in the table order, by
\[
 T_j^2+c-y_j=0,\qquad N(T_j)-(T_j-b)w_{y_j}=0.
\]
The minor on columns zero through twenty-four has determinant $21$
modulo $97$.  Thus the lifting conclusion has a second certificate using
ordinary polynomial evaluations rather than the norm equations.

Fix $v_2=3,v_3=70$ as integers in
\eqref{eq:six-word-norm-system}.  Multivariate Hensel lifting gives a
solution of the remaining thirteen equations in $\mathbb Z_{97}$, with
invertible Jacobian.  Its coordinates are algebraic over $\mathbb Q$:
in the field they generate, differentiating the thirteen equations and
using the invertible Jacobian makes every coordinate differential zero;
in characteristic zero this forces transcendence degree zero.
They therefore lie in a number field embedded in $\mathbb Q_{97}$.

All required guards hold at the residue point and hence remain nonzero
in this lift.  They include distinct nonzero seed parameters and signed
nodes, $e_3\ne0$, the first pole and critical-value exclusions,
$O(y)-w_y\ne0$, the new pole being off the whole core, and $N(b)\ne0$.
The five forced roots of $N-(T-b)G_5$ persist.  Its remaining quadratic
still has nonzero leading coefficient and discriminant and has no root
in the core or at $b$: these are polynomial nonvanishing conditions,
verified by the displayed finite-field factorization.
Adjoin its two roots and the other core square roots to the number field.
They supply two distinct fresh coordinates.  Thus the five polynomials
$G_i$ each have at least fourteen agreements on a twenty-eight-point word,
and the proper function $N/(T-b)$ also has at least fourteen.

\paragraph{The second pole-clearing extension.}
Choose a quadratic $\psi(U)=U^2+d$ whose critical value avoids the
length-twenty-eight domain and $b$, and split its finitely many required
fibers in a number-field extension.  Use the fifty-six preimages of that
domain, the two preimages of $b$, and two further fresh points.  The six
polynomials are
\[
 (\psi-b)G_i(\psi)\quad(1\le i\le5),\qquad N(\psi).
\]
They have degree at most fourteen.  On the first fifty-six points set
the received word to $(\psi-b)$ times the preceding word, on the pole
fiber set it to zero, and on the two fresh points give it the value
$N(\psi)$.  Every old candidate has at least $28+2$ agreements; the new
one has at least $28+2$.  The old candidates remain distinct, and equality
of the new candidate with any old one would make $N/(T-b)$ polynomial,
contrary to properness.

Finally, all data lie in one number field.  Exclude the finitely many
primes where denominators, distinctness, or any used nonzero guard fails.
At every remaining completely splitting prime the data reduce into the
prime field and preserve the sixty nodes, six distinct polynomials,
degree bounds, and displayed agreements.  There are arbitrarily large
such primes, proving the theorem.
\end{proof}
